\input{../tex-inputs/latex-leseansicht-vorspann}

\section[Gustav Schwarzkopf an Arthur Schnitzler, 22. 7. 1898]{L04289 Gustav Schwarzkopf an Arthur Schnitzler, 22. 7. 1898}
\nopagebreak\mylabel{L04289v}
\rehead{ }\normalsize\beginnumbering\briefempfaengerindex{Schnitzler, Arthur@\textsc{Schnitzler, Arthur}!zzzSchwarzkopf, Gustav@\emph{von Gustav Schwarzkopf}!1898-07-221@{22. 7. 1898}|(be}
\toendnotes[C]{\smallbreak\pagebreak[2]}
\correspDesc{Versand  durch Gustav Schwarzkopf am 22. 7. 1898 in Wien
\newline{}Übermittlung  am 23. 7. 1898 in Wien
\newline{}Zustellung  am 24. 7. 1898 in Bad Gastein
\newline{}Erhalt  durch Arthur Schnitzler am 24. 7. 1898 in Bad Gastein}\toendnotes[C]{\smallbreak}
\Standort{CUL, Schnitzler, B 96a.}
\physDesc{Postkarte, 608 Zeichen
\newline{}Handschrift: schwarze Tinte, deutsche Kurrent
\newline{}Versand: 1) Stempel: »\nobreak{}\oindex{I., Innere Stadt@\textbf{I., Innere Stadt}, \emph{Verwaltungsgebiet}|pwk}Wien 1/1 1, 23. 7. 98, 1–2N\nobreak{}«.   2) Stempel: »\nobreak{}\oindex{Bad Gastein@\textbf{Bad Gastein}, \emph{Hauptstadt}|pwk}Badgastein, 24/7 98, 10–{[}1{]}8\nobreak{}«. }\toendnotes[C]{\smallbreak}\pstart{}{\pb}Herrn
                  \textsc{D\textsuperscript{r} Arthur Schnitzler}\pend{}\pstart{}Wildbad \textsc{Gastein}\oindex{Bad Gastein@\textbf{Bad Gastein}, \emph{Hauptstadt}|pw}\pend{}\pstart{}\textsc{Villa Wassing\oindex{Villa Dr. Wassing@\textbf{Villa Dr. Wassing}, \emph{Sanatorium}|pw}}.\pend{}{\bigskip}\vspace{1em}
\pstart
           \raggedleft{}{\pb}22. 7. 98\pend
           \vspace{0.5em}
\pstart
           Lieber Arthur, ich werde alſo entweder am 26\textsuperscript{ten} Nachmittag oder am
      27\textsuperscript{ten} Morgen in Salzburg\oindex{Salzburg@\textbf{Salzburg}, \emph{Verwaltungsgebiet}|pw}{ }ſein, ka{\geminationn} aber nur bis 28. Abends
      bleiben, da \textsc{Hirschfelds\pwindex{Hirschfeld, Robert 17.\,9.\,1857 Žďár nad Sázavou – 2.\,4.\,1914 Salzburg@\textsc{Hirschfeld, Robert} (17.\,9.\,1857 Žďár nad Sázavou – 2.\,4.\,1914 Salzburg), \emph{Journalist, Musikkritiker}|pw}\pwindex{Hirschfeld, Emilie 12.\,12.\,1868 Berlin – 3.\,4.\,1942 Wien@\textsc{Hirschfeld, Emilie} (12.\,12.\,1868 Berlin – 3.\,4.\,1942 Wien), \emph{Sängerin}|pw}}, die in den erſten Tagen des Auguſt nach \label{K_L04289-1v}\edtext{Heindorf\oindex{Heindorf@\textbf{Heindorf}|pw}}{\lemma{\textnormal{\emph{Heindorf}}}\Cendnote{\textnormal{Welcher Ort gemeint ist,
      ist unklar. Gegen eine weitere Reise spricht, dass man sich in diesem Fall eher an der Eisenbahn getroffen hätte. In unmittelbarer
      Nähe des Traunsees\oindex{Traunsee@\textbf{Traunsee}, \emph{See}|pwk} bietet sich kein Ort an.}}}\label{K_L04289-1} gehen, mich{ }ſchon am 26. in Traunkirchen\oindex{Traunkirchen@\textbf{Traunkirchen}|pw} erwarten.\pend
           
\pstart
           Zu der »Billigkeit« der »Elektr. \textsc{Hôtel}\oindex{Hotel Bristol Salzburg@\textbf{Hotel Bristol Salzburg}, \emph{Hotel}|pw}« habe ich wenig
               Zutrauen. Man hat mir ein \textsc{Hôtel} »Zur neuen Stadt\oindex{Johann Pachler’s Hotel zur Neuen Stadt@\textbf{Johann Pachler’s Hotel zur Neuen Stadt}, \emph{Hotel}|pw}« gleich
               beim Bahnhof\oindex{Hauptbahnhof Salzburg@\textbf{Hauptbahnhof Salzburg}, \emph{Bahnhofsgebäude}|pwv}, Haydnſtraße\oindex{Haydnstraße [Salzburg]@\textbf{Haydnstraße [Salzburg]}, \emph{Straße}|pw}, als gut empfohlen. We{\geminationn} ich
      vor Ihnen am 26. anko{\geminationm}e, werde ich Sie da{\geminationn} gegen Abend in
      Ihrem \textsc{Hôtel\oindex{Hotel Bristol Salzburg@\textbf{Hotel Bristol Salzburg}, \emph{Hotel}|pwv}} erwarten. Auf Wiederſehen.\pend
           
\pstart
           Herzlichſt Ihr{\\[\baselineskip]}\spacefill\mbox{GustavSch}\pend
           \leftskip=0em{}
\pstart
           \noindent{}Mein beſten Empfehlungen Ihrer Frau Mama\pwindex{Schnitzler, Louise 8.\,7.\,1840 Kőszeg – 9.\,9.\,1911 Wien@\textsc{Schnitzler, Louise} (8.\,7.\,1840 Kőszeg – 9.\,9.\,1911 Wien)|pwv}.\pend
           \selectlanguage{ngerman}\endnumbering\briefempfaengerindex{Schnitzler, Arthur@\textsc{Schnitzler, Arthur}!zzzSchwarzkopf, Gustav@\emph{von Gustav Schwarzkopf}!1898-07-221@{22. 7. 1898}|)be}\mylabel{L04289h}
\begin{anhang}
\end{anhang}\newcommand{\dateiname}{L04289}\newcommand{\titel}{Gustav Schwarzkopf an Arthur Schnitzler, 22. 7. 1898}\newcommand{\editorInnen}{Herausgegeben von Jahnke, SelmaMüller, Martin Anton}\input{../tex-inputs/latex-leseansicht-abspann}
